Ce projet ITInéraire nous a permis de mettre en pratique des compétences variées autour des systèmes embarqués, de l’acquisition de données capteurs, et de l’analyse statistique, tout en développant un outil fonctionnel de suivi d’activité sportive. En nous appuyant principalement sur un GPS et une IMU, nous avons pu concevoir un système capable d’enregistrer et d’analyser des trajets piétons de manière fiable.\\

Nous avons choisi de ne pas intégrer le capteur LiDAR dans la version finale, préférant concentrer nos efforts sur la synchronisation et la fusion GPS/IMU, qui représentaient déjà un vrai défi technique. Cette décision nous a permis de garantir une meilleure stabilité du pipeline de traitement et de concentrer notre temps sur l’amélioration de la qualité des données et de l’interface utilisateur.\\

Le traitement automatisé mis en place (nettoyage, filtrage, calcul d’indicateurs) couplé à un menu interactif en ligne de commande, rend le système à la fois rigoureux et accessible. L’utilisateur peut ainsi visualiser ses performances sans avoir à manipuler les données brutes, ce qui constitue un vrai atout en termes de confort d’utilisation.\\

Ce projet nous a également appris à mieux collaborer en équipe, chacun ayant contribué à des aspects complémentaires : matériel, traitement des données, interface, ou encore rédaction du rapport. Cette répartition nous a permis d'avancer efficacement tout en apprenant les uns des autres.\\

En résumé, ce projet a été formateur à la fois sur le plan technique et organisationnel. Il nous a permis de réaliser une solution concrète et complète de suivi de trajectoire, tout en développant des méthodes d’analyse reproductibles et robustes, qui pourraient être réutilisées ou enrichies dans des projets futurs.\\

